\begin{table}
\centering
\caption{\label{tab:main_cand_char}Candidate Characteristics: Quota vs. Open Seats (Rajasthan 2020)}
\centering
\fontsize{8}{10}\selectfont
\begin{threeparttable}
\begin{tabular}[t]{lllllll}
\toprule
\multicolumn{1}{c}{ } & \multicolumn{3}{c}{Winners} & \multicolumn{3}{c}{Candidates} \\
\cmidrule(l{3pt}r{3pt}){2-4} \cmidrule(l{3pt}r{3pt}){5-7}
Characteristic & Open & Quota & Difference & Open & Quota & Difference\\
\midrule
Age & 42.8 & 41.3 & 1.6 & 41.5 & 39.9 & 1.7\\
Graduate (\%) & 23.9 & 12.7 & 11.2 & 22.2 & 13.3 & 8.9\\
Unemployed (\%) & 9.9 & 60.9 & -50.9 & 10.7 & 61.4 & -50.7\\
Log assets & 12.9 & 12.6 & 0.4 & 12.3 & 12.1 & 0.3\\
\bottomrule
\end{tabular}
\begin{tablenotes}
\item \textit{Note: } 
\item Descriptive open-minus-quota comparisons, conditional on candidacy or winning. Graduation and unemployment are percentages; their differences are percentage points. Other rows use their stated units. Graduation rates exclude records with missing education. Age means exclude recorded ages outside 21--110 years; other outcomes retain those records. Assets winsorized at the 10th and 90th percentiles of the candidate sample before taking logs. These comparisons do not identify a causal effect on candidate quality.
\end{tablenotes}
\end{threeparttable}
\end{table}
